\documentclass[11pt]{article}
\usepackage{amsmath,amssymb,amsthm}
\newtheorem{theorem}{Theorem}
\newtheorem{lemma}[theorem]{Lemma}
\title{Independent Permutation Fibers Inside the Original Boolean Cube}
\author{}
\date{November 23, 2026}
\begin{document}
\maketitle

\paragraph{Status and scope.}
This note gives a complete written construction of many isolated regular
interpolation roots in a strictly admissible family. It is not Lean
formalized. The interpolation constraints use a specified subset of
Boolean vertices; the roots need not agree at the other vertices.
Consequently neither a universal sign compiler nor the requested global
lower bound follows. The final section identifies the remaining readout
problem explicitly.

\begin{lemma}[An elementary reciprocal-power interpolation fact]
Let $\lambda_1<\cdots<\lambda_H$, let $k$ be a positive integer, and
let $d_1,\ldots,d_H$ be distinct real numbers greater than $\lambda_H$.
The matrix $(d_i-\lambda_j)^{-k}$ is nonsingular.
\end{lemma}
\begin{proof}
More generally a nonzero linear combination of the $H$ functions
$(t-\lambda_j)^{-k}$ has at most $H-1$ distinct zeros for
$t>\lambda_H$. Prove this simultaneously for every $k\geq1$ by
induction on $H$. The case $H=1$ is immediate.
If
$g(t)=\sum_j a_j(t-\lambda_j)^{-k}$
had $H$ distinct zeros, multiply it by $(t-\lambda_H)^k$ and
differentiate. Rolle's theorem supplies at least $H-1$ distinct zeros
of
\[
 k(t-\lambda_H)^{k-1}
 \sum_{j<H}a_j(\lambda_H-\lambda_j)(t-\lambda_j)^{-k-1}.
\]
The prefactor is positive. Unless all $a_j$ for $j<H$ vanish, the
inductive assertion bounds these zeros by $H-2$, a contradiction.
If they all vanish, the original nonzero function is a single
reciprocal power and has no zero. Thus the assertion holds, and a
matrix kernel vector would contradict it. This proof avoids any
external Cauchy-square determinant identity.
\end{proof}

Fix integers $d,r\geq1$ and put
\[
 H=2^d,\qquad n=r+d+1,\qquad
 C=(r+3)(H+2),\qquad \lambda_h=h+1\quad(0\leq h<H).
\]
Write the actual $n$ input coordinates as
$x\in\{0,1\}^r$, $y\in\{0,1\}^d$, and $z\in\{0,1\}$.
Define the affine address denominator
\[
 D(y)=C-\sum_{k=0}^{d-1}2^k y_k.
\]
Its $H$ values are distinct. For real parameters
$a_{hi},u_{hi}$ define the $2H$-head score
\[
 F_{u,a}(x,y,z)=
 \sum_{h=0}^{H-1}
 \frac{\sum_{i=1}^r a_{hi}x_i+z}
      {D(y)-\sum_{i=1}^r u_{hi}x_i-z}
 -
 \sum_{h=0}^{H-1}
 \frac{z}{D(y)-\lambda_h-\sum_{i=1}^r x_i-z}.
 \tag{1}
\]
The second sum consists of fixed reference heads. The first sum has
exactly $p=2Hr$ free real parameters.

\begin{theorem}[Factorially many independent positive-chamber roots]
Evaluate (1) at the $p$ distinct vertices
\[
 \mathcal T=\{(e_i,y,z):1\leq i\leq r,\ y\in\{0,1\}^d,\
                                               z\in\{0,1\}\}.
 \tag{2}
\]
This evaluation map from its $p$ free parameters to $\mathbb R^p$
has at least $(H!)^r$ distinct regular zeros. Every denominator at
every listed zero has strictly negative slopes and is positive on
the entire $n$-dimensional Boolean cube. All these roots persist as
distinct regular real branches for every target vector in one common
open neighborhood of zero.
\end{theorem}
\begin{proof}
Choose an independent permutation $\pi_i$ of the $H$ tokens for each
$i$ and set
\[
             a_{hi}=0,\qquad u_{hi}=\lambda_{\pi_i(h)}+1.
 \tag{3}
\]
At $z=0$ both numerators vanish. At $x=e_i,z=1$ the first denominator
is $D(y)-\lambda_{\pi_i(h)}-2$, while the corresponding reference
denominators are $D(y)-\lambda_h-2$. Their reciprocal sums are equal
because $\pi_i$ is a permutation. Thus (3) gives a zero of every
evaluation in (2), and the $(H!)^r$ parameter vectors are distinct.

For the first family of heads all coordinate slopes are strictly
negative: they are $-u_{hi}$, $-2^k$, and $-1$. Their cube minima
are at least
\[
 C-(H-1)-r(H+1)-1=2H+r+6>0.
 \tag{4}
\]
The reference head slopes are $-1$, $-2^k$, and $-1$, and their
cube minima are at least
\[
 C-(H-1)-H-r-1=(r+1)H+r+6>0.
 \tag{5}
\]
Thus these are actual original-model admissible affine ratios,
with no weak slopes, extra coordinates, or limiting transfer.
The strict conditions persist in open parameter neighborhoods.

At samples with $x=e_i$, the derivatives with respect to
$a_{hj},u_{hj}$ vanish when $j\ne i$. The complete Jacobian is
therefore block diagonal in $i$. For the $i$th block order the
parameters as $(u_{\cdot i},a_{\cdot i})$ and the samples first by
$z=0$, then by $z=1$. At (3) the block is
\[
 \begin{pmatrix}
  0 & C_i^0\\
  C_i^1\circ C_i^1 & C_i^1
 \end{pmatrix},
 \qquad
 (C_i^0)_{yh}=\frac1{D(y)-\lambda_{\pi_i(h)}-1},
 \quad
 (C_i^1)_{yh}=\frac1{D(y)-\lambda_{\pi_i(h)}-2}.
 \tag{6}
\]
The first lemma makes both $C_i^0$ and $C_i^1\circ C_i^1$
nonsingular. All address poles exceed the largest shifted token,
as is also immediate from (4). The determinant in (6) is
$(-1)^H\det(C_i^0)\det(C_i^1\circ C_i^1)$ and is nonzero.
Thus every full Jacobian is nonsingular.

Apply the inverse function theorem in pairwise disjoint neighborhoods
of the finitely many roots, keeping the strict margins (4)--(5).
Every branch image contains an open neighborhood of zero. Their
finite intersection contains one common open neighborhood, proving
the last assertion.
\end{proof}

\paragraph{Quantitative relevance to a proposed counting shortcut.}
Take $r=\lceil\log_2 d\rceil$ and let $d$ tend to infinity. The
actual head count is $h=2H$, and
\[
 N=2^n=2^{r+1}H,\qquad
 \frac{h}{N/n}=\frac{n}{2^r}=\Theta(1).
 \tag{7}
\]
On the other hand,
\[
 \log\bigl((H!)^r\bigr)
 \geq \frac{Hr}{2}\log(H/2)
 =\Omega(Hn\log n).
 \tag{8}
\]
Here and below the unlabeled logarithm is natural. For any fixed
$K>0$,
\[
 \log\bigl(h!\exp(Khn)\bigr)
 \leq h\log h+Khn=O(Hn),
 \tag{9}
\]
so (8) eventually exceeds (9). Even quotienting by simultaneous
permutation of all main heads leaves $(H!)^{r-1}$ roots and the
same asymptotic comparison.

Thus a uniform bound $h!\exp(O(hn))$ for regular interpolation
fibers of all affine-parameter subfamilies cannot follow from
positivity on the full Boolean cube alone. The example already uses
actual cube evaluations and strict original admissibility. It has
only $2Hr$ free parameters, with the other head coefficients fixed;
it does not claim the same root count for a square system in every
unrestricted head coefficient. Nor does a root count alone determine
the number of strict truth-table sign patterns.

\paragraph{The mixed-coordinate readout problem.}
At a root, the score is zero for every $z=0$ input. For $z=1$ and
$S=\{i:x_i=1\}$, it is exactly
\[
 \sum_h\frac1{D(y)-|S|-1-\sum_{i\in S}\lambda_{\pi_i(h)}}
 -
 \sum_h\frac1{D(y)-|S|-1-\lambda_h}.
 \tag{10}
\]
This depends on the multiset of sums of the independently permuted
tokens. For example, with two signals and two tokens, identical
permutations give the sums $\{2,4\}$ and opposite permutations give
$\{3,3\}$; strict convexity of the reciprocal function distinguishes
their values. So the unused cube vertices can detect some joint
permutation information.

However these unmodified root scores do not read independent signs.
When $S=\varnothing$, every term of the first sum is smaller than the
corresponding reference term. When $|S|=1$, permutation invariance
gives zero. When $|S|\geq2$, align the first sum's terms with one
chosen permutation $\pi_{i_0}$. Since every other token is at least
one,
\[
 \sum_{i\in S}\lambda_{\pi_i(h)}
 \geq \lambda_{\pi_{i_0}(h)}+|S|-1.
 \tag{11}
\]
For $|S|\geq2$ this implies that every term of the first reciprocal
sum is strictly greater than its aligned reference term.
Consequently all the strict mixed-face signs in (10) are positive,
independently of the permutation choices. Arbitrary labels at these
vertices have not been obtained. A usable global construction would
require new readouts or additional cancellations that preserve a
superconstant saving; the many regular roots alone do not supply them.

\end{document}
